\section{Introduction}
\label{sec:introduction}

Dough manipulation combines large deformation, changing tool contact, incomplete observation, and a response that depends on preparation and loading history. These properties complicate the use of simulation for planning, control, and robot learning. A useful model must predict how an object will deform under a proposed action, but its initial geometry and mechanical parameters are rarely known with the accuracy assumed in a synthetic experiment. Deformable-object modeling therefore requires a joint consideration of numerical methods, sensing, and identification, rather than only a more detailed forward simulator~\cite{review2020}.

For flour dough, this difficulty is particularly important. Elastic recovery, rate-dependent resistance, lasting deformation, and adhesion may occur within the same manipulation. Their relative importance varies with formulation, hydration, mixing, and previous handling~\cite{yazar2023}. A compact constitutive approximation can still be useful within a restricted operating range. However, minimizing an image discrepancy does not by itself establish that its parameters are intrinsic material properties. An apparent change in stiffness or viscosity can also compensate for an inaccurate initial thickness, incorrect tool pose, or excessive numerical dissipation.

Recent work has established several parts of this problem. RoboCraft and RoboCook demonstrate learned predictive models for elastoplastic manipulation and real dough tasks~\cite{robocraft2022,robocook2023}. Differentiable physics-based system identification (DPSI) estimates elastoplastic material parameters from robot-conditioned observations of plasticine~\cite{dpsi2025}. EMPM includes visually calibrated material point simulation and bimanual bread-dough interaction~\cite{empm2026}. Single-depth-camera calibration is also established for soft solids with known reference geometry~\cite{diffcal2023}. Consequently, the combination of depth sensing, two tools, and material point simulation is not sufficient as a novelty claim. The relevant question concerns the information provided by a measurement protocol and the predictive value of the resulting model.

We investigate the following hypothesis: \emph{a simplified elastoplastic model, calibrated from time-resolved observations of a single depth view and recorded tool trajectories, can improve prediction of previously unseen dough manipulation without further material-parameter adjustment.} Recorded tool motion supplies the prescribed mechanical input; the method is not a camera-only reconstruction of unknown actions. The observation sequence supplies intermediate deformation information, potentially distinguishing parameter combinations that produce similar endpoints. Whether it actually does so must be tested against endpoint-only fitting and under different loading and recovery conditions.

TaichiDough implements this identification problem using moving least squares material point simulation, a partial-visible-geometry objective, and checkpointed reverse-mode differentiation. Three methodological elements organize the study. First, scene alignment, tool registration, and hidden-volume reconstruction are made explicit and separated from material optimization. Second, a bounded five-parameter model is fitted using a loss that accounts for incomplete depth support while penalizing missing predictions in observed regions. Third, shared-material fitting and frozen-parameter evaluation distinguish agreement on calibration recordings from prediction on independent episodes.

The intended contribution is an experimentally testable approach to measurement-conditioned dough modeling, not a new continuum discretization or proof of unique rheological identification. The present draft specifies the implemented method and the evaluation needed to support the hypothesis. It does not report completed real-data generalization results. Parameters are therefore described as \emph{effective simulation parameters}, conditional on the constitutive law, discretization, geometry, and contact model~\cite{discrepancy2014}.
